% \vspace{-5mm}

\section{Introduction}
\label{sec:intro}

Blockchain technology enables the creation of decentralized, resilient, and
programmable ledgers on a global scale. Smart contracts, which can be deployed
onto a blockchain, allow developers to encode intricate transaction rules that
govern the ledger. These features have made blockchains and smart contracts
essential infrastructure for a variety of decentralized financial services
(DeFi). As of April 1st, 2023, the Total Value Locked (TVL) in 1,417 DeFi smart
contracts had reached $50.15$ billion~\cite{DeFillama}.

However, security attacks are critical threats to smart contracts. Attackers
can exploit vulnerabilities in smart contracts by sending malicious
transactions, potentially stealing millions of dollars from users.
Particularly, a new type of security threat has emerged, exploiting design
flaws in DeFi contracts by leveraging large amounts of digital assets. These
attacks, commonly referred to as \emph{flash loan attacks}~\cite{SlowMistStats,
mckay2022defi, zhang2023demystifying}, 
typically involve borrowing the required large amount of assets
from flash loan contracts. Among the top $200$ costliest attacks recorded in
Rekt Database,  the financial loss caused by $36$ flash loan attacks exceeds
$418$ million USD~\cite{mckay2022defi}. 

%The flash loan is a special kind of loan that occurs inside one transaction,
%i.e., the borrower must return all borrowed assets back before the end of the
%transaction otherwise the flash loan protocol will abort the transaction to
%cancel the loan.

A typical flash loan attack transaction consists of a sequence of actions, or
function calls to smart contracts. The first action involves borrowing a
substantial amount of digital assets from a flash loan contract, while the last
action returns these borrowed assets. The sequence of actions in the middle
interacts with multiple DeFi contracts, using the borrowed assets to exploit
their design flaws. When a DeFi contract fails to consider corner cases created
by the large volume of the borrowed assets, the attacker may extract
prohibitive profits. For example, many flash loan attacks use borrowed assets
to temporarily manipulate asset prices in a DeFi contract to trick the contract
to make unfavorable trades with the
attacker~\cite{qin2021attacking,cao2021flashot}. Although researchers have
developed many automated program analysis and verification
techniques~\cite{mossberg2019manticore,feng2019precise, grieco2020echidna,
brent2020ethainter} to detect and eliminate bugs in smart contracts, these
techniques cannot handle flash loan attack vulnerabilities. This is because
such vulnerabilities are design flaws rather than implementation bugs.
Moreover, these techniques typically operate with one contract at a time, but
flash loan attacks almost always involve multiple DeFi contracts interacting
with each other.

\noindent \textbf{\name:}
We present {\name}, the first automated end-to-end program synthesis tool for
detecting flash loan attack vulnerabilities. Given a set of smart contracts and
candidate actions in these contracts, {\name} automatically synthesizes an
action sequence along with all action parameters to interact with the contracts
to exploit potential flash loan vulnerabilities. 
Additionally, {\name} can analyze past blockchain transaction history,
assisting users in identifying candidate actions for synthesis.

The primary challenge {\name} faces is that the underlying logic of DeFi actions
is often too sophisticated for standard solvers to handle. Even
if the action sequence was already known, a naive application of symbolic
execution might not be able to find action parameters because it may need to extract
overly complicated symbolic constraints causing the solvers to
time out. Moreover, {\name} synthesizes the action sequence and the action
parameters together and therefore faces an additional search space explosion
challenge.

{\name} addresses these challenges with its novel
\emph{synthesis-via-approximation} technique. Instead of attempting to extract
accurate symbolic expressions from smart contract code, {\name} collects data
points to approximate the effect of contract functions with numerical methods. {\name} then uses the approximated expressions to drive the
synthesis. {\name} also incrementally improves the approximation with its novel
\emph{counterexample driven approximation refinement} techniques, i.e., if the
synthesis fails because of a large deviation caused by the approximations,
{\name} collects the corresponding data points as counterexamples to
iteratively refine the approximations. The combination of these techniques
allows the underlying optimizer of {\name} to work with more tractable
expressions. It also decouples the two difficult tasks, finding the action
sequence and finding the action parameters. When working with a set of
coarse-grained approximated expressions, {\name} can filter out unproductive
action sequences with a small cost.

\noindent \textbf{Experimental Results:} 
We evaluate {\name} on $16$ DeFi benchmark protocols that were victims to flash
loan attacks and $2$ DeFi benchmark protocols from Damn Vulnerable DeFi
challenges~\cite{DVD}. {\name} synthesizes adversarial attacks for $16$ out
of the $18$ benchmarks. For comparison, a baseline with manually crafted
accurate action summaries only synthesizes attacks for $7$ out of the $18$.

\noindent \textbf{Contributions:}
This paper makes the following contributions:
\begin{itemize}[leftmargin=*]
\item \textbf{\name:} The first automated
end-to-end program synthesis tool for detecting flash loan attack
vulnerabilities. It enables approximate attack synthesis 
without diving into sophisticated logics of DeFi contracts. 
\item \textbf{Synthesis-via-approximation:} A novel \revise{synthesis-via-approxi- mation} technique to handle sophisticated logics of DeFi contracts. 
\item \textbf{Counterexample Driven Approximation Refinement:} A novel counterexample driven approximation refinement technique to
        incrementally improve the approximation during the synthesis process.
\item \textbf{Experimental Evaluation:} 
We implemented {\name} in a tool and evaluated it on $16$ protocols that were victims to flash
loan attacks and $2$ fictional flash loan attacks. 
\end{itemize}

Our solution, {\name} has been adopted and further developed by Quantstamp, a leading smart contract auditing company for the detection of flash loan vulnerabilities in DeFi contracts~\cite{adoption,adoptionNews1,adoptionNews2}. 
% The tool is currently undergoing further development as the company works towards enhancing its features and robustness.